\subsubsection{Wissenschaftliche Gütekriterien}
\label{subsubsec:Guetekriterien}

Das Evaluationsdesign wurde unter Berücksichtigung der wissenschaftlichen Gütekriterien Validität, Reliabilität und Objektivität entwickelt.

\paragraph{Validität}

Validität bezeichnet das Ausmaß, in dem ein Messinstrument tatsächlich das misst, was es messen soll. Das Drei-Säulen-Modell wurde so konzipiert, dass es die zentralen Dimensionen der Konfigurationseffizienz abdeckt:

\begin{itemize}
    \item \textbf{Inhaltsvalidität:} Die 6-Punkte-Checkliste basiert auf empirisch beobachteten Fehlerquellen aus der Demo-Werk-Entwicklung und deckt kritische Validierungsaspekte ab (Vollständigkeit, ID-Konsistenz, Zeiteinheiten, BOM-Berechnung, Material-Referenzen, ID-Duplikate).

    \item \textbf{Konstruktvalidität:} Die drei Säulen (Zeit, Vollständigkeit/Fehlerrate, Reproduzierbarkeit) operationalisieren das theoretische Konstrukt ``Konfigurationseffizienz'' entlang wissenschaftlich fundierter Dimensionen (Zeitwirtschaft nach REFA, Qualitätsmanagement, Wissensmanagement).

    \item \textbf{Externe Validität:} Praxisvalidierung in zwei realen Kundenprojekten bestätigt die Anwendbarkeit der Artefakte außerhalb des kontrollierten Experiment-Settings.
\end{itemize}

\paragraph{Reliabilität}

Reliabilität bezeichnet die Zuverlässigkeit und Wiederholbarkeit von Messungen. Das Evaluationsdesign adressiert Reliabilität durch mehrere Mechanismen:

\begin{itemize}
    \item \textbf{Test-Retest-Reliabilität:} Jeder Proband absolvierte 6 Durchläufe pro Bedingung, was die Konsistenz der Messungen über mehrere Zeitpunkte hinweg ermöglicht.

    \item \textbf{Interrater-Reliabilität:} Reproduzierbarkeitstest durch Proband 2 (Kollege ohne vorherige Master-Creator-Erfahrung) validiert, dass die entwickelten Artefakte (SOPs) von verschiedenen Personen konsistent angewendet werden können.

    \item \textbf{Methodische Kontrolle:} Counterbalanced Design reduziert Reihenfolgeeffekte; standardisiertes Experiment-Setup (5/5/1/1/1/1) gewährleistet vergleichbare Bedingungen.
\end{itemize}

\paragraph{Objektivität}

Objektivität bezeichnet die Unabhängigkeit der Messung vom Forscher. Das Evaluationsdesign nutzt mehrere Strategien zur Sicherstellung der Objektivität:

\begin{itemize}
    \item \textbf{Durchführungsobjektivität:} Standardisiertes Experiment-Setup mit klaren Anweisungen und identischen Aufgabenstellungen für alle Probanden.

    \item \textbf{Auswertungsobjektivität:} Objektive Metriken (Zeit in MM:SS,ms, Git-Commit-Anzahl, Checklisten-Punkte) minimieren Interpretationsspielräume.

    \item \textbf{Interpretationsobjektivität:} Externe Validierung durch Kollegen (Proband 2) reduziert Bestätigungstendenz durch den Entwickler (Proband 1). Die 6-Punkte-Checkliste definiert binäre Validierungskriterien (erfüllt/nicht erfüllt) ohne subjektive Bewertungsskalen.
\end{itemize}

Die Kombination dieser Gütekriterien stellt sicher, dass die Evaluationsergebnisse wissenschaftlich fundiert und nachvollziehbar sind.
